%!TeX program = lualatex
\renewcommand{\thelektion}{L04 \textperiodcentered{} Bits, Bytes \& Hexadezimal}

\begin{frame}\lessontitle{Lektion 04}{Bits, Bytes \& Hexadezimalsystem}{Die kleinste Informationseinheit, Byte-Strukturen und Basis 16 \textperiodcentered{} Skript, Kapitel 2}\end{frame}

{\setbeamertemplate{footline}{\darkfootline}
\begin{frame}\sectionframe[FDUblue]{Bits \& Bytes}{Wie zählt und speichert man mit Nullen und Einsen?}
\end{frame}}

\begin{frame}\FDUtitle{Bits \& Bytes}
\begin{itemize}[<+->]
	\item \textbf{bit} = \enquote{\textit{\underline{bi}nary digi\underline{t}}}, \enquote{binäre Ziffer} (eine 0 oder eine 1)
	\item \textbf{byte} = \enquote{Bissen} (von engl.\ \textit{bite} = Bissen, aus 8 bits) \emoji{yum}
	\item Mit einem Byte lassen sich $2^8=256$ verschiedene Werte darstellen
\end{itemize}
\end{frame}

\begin{frame}\FDUtitle{Führende Nullen}
Die binäre Repräsentation der Dezimalzahl 22:
\[00010110_2 = 22_{10}\]
\begin{center}
Bei Binärzahlen werden -- anders als bei Dezimalzahlen -- üblicherweise \textbf{alle} Bits ausgeschrieben.
\end{center}
\vspace{3mm}
\begin{center}
\textbf{Frage}: Wie stellt man $42_{10}$ mit einem Byte dar?
\end{center}
\only<2->{\begin{center}\Large $42_{10} = 00101010_2$\end{center}}
\end{frame}

\begin{frame}\FDUtitle{Grössenordnungen}
\begin{table}[H]
	\footnotesize
	\centering
	\begin{tabular}{@{}l|rl@{}}
		\toprule
		\textbf{Einheit} & \textbf{Grösse} & Beispiel \\ \midrule
		KB (Kilobyte) & $10^3$ B & eine Text-Datei \\
		MB (Megabyte) & $10^6$ B & eine Musik-Datei \\
		GB (Gigabyte) & $10^9$ B & eine Video-Datei \\
		TB (Terabyte) & $10^{12}$ B & kleiner Firmen-Server \\
		\bottomrule
	\end{tabular}
\end{table}
\end{frame}

{\setbeamertemplate{footline}{\darkfootline}
\begin{frame}\sectionframe[FDUgreen]{Das Hexadezimalsystem}{Warum brauchen wir die Basis 16 in der Informatik?}
\end{frame}}

\begin{frame}\FDUtitle{Das Hexadezimalsystem ($b=16$)}
Binärzahlen wie $11111111_2$ sind lang und unübersichtlich. 
\vspace{2mm}
Das Hexadezimalsystem nutzt 16 Ziffern:
$$\lrc{0,1,2,3,4,5,6,7,8,9,A,B,C,D,E,F}$$
\begin{center}
$A=10, B=11, C=12, D=13, E=14, F=15$
\end{center}
\end{frame}

\begin{frame}\FDUtitle{Notwendigkeit \& Vorteil von Hexadezimal}
Da $16 = 2^4$, entspricht \textbf{1 Hex-Ziffer} genau \textbf{4 Bits} (einem Halbbyte/Nibble).
\vspace{3mm}
\begin{center}
\textbf{1 Byte (8 Bits) = exakt 2 Hex-Ziffern!}
\end{center}
\begin{align*}
	11111111_2 &= FF_{16} = 255_{10} \\
	10101100_2 &= AC_{16} = 172_{10}
\end{align*}
\begin{center}
Viel kompakter, besser lesbar und weniger anfällig für Tippfehler!
\end{center}
\end{frame}

\begin{frame}{Auftrag}{Kapitel 2: Bits, Bytes \& Hexadezimal}
\aufgabebox{\textbf{Aufgabe 2.1 \& 2.2}: $42_{10}$ als Byte ausdrücken \& grösste Zahl mit 1 Byte}\\[3mm]
\aufgabebox{\textbf{Aufgabe 2.3 \& 2.4}: Hexadezimale Umrechnungen ($27_{10}$, 44, 23, 90, 124, 306)}
\end{frame}


